%=============================================================================
% APPENDIX M: Double-Transit Enhancement Analysis
%=============================================================================

\section{Double-Transit Enhancement: Derivation and Tests}\label{app:double-transit}

This appendix derives the $\Gamma = 4$ double-transit enhancement factor from two physical inputs: (i) resonantly scattered photons sample the $\psi$-gradient on both the incoming and outgoing legs, acquiring twice the frequency detuning of a locally emitted line, and (ii) the asymmetry observable is quadratic in the effective detuning.  The derivation is presented with explicit assumptions and falsifiers.

%-----------------------------------------------------------------------------
\subsection{Definitions and Setup}
%-----------------------------------------------------------------------------

Let $\psi(\mathbf{x})$ be the DFD scalar field with refractive index $n = e^{\psi}$ 
and one-way light speed $c_1 = c\,e^{-\psi}$. Consider two UV lines observed by UVCS:
\begin{itemize}[nosep]
\item \textbf{H Ly-$\alpha$:} Dominated by resonant scattering of chromospheric 
radiation in the corona.
\item \textbf{O~VI:} Dominated by local (collisional) emission in the corona.
\end{itemize}

Let $A$ denote the measured asymmetry amplitude statistic, and define:
\begin{equation}
R \equiv \frac{A_{\rm Ly\alpha}}{A_{\rm OVI}} = \Gamma \left(\frac{\sigma_{\rm OVI}}{\sigma_{\rm Ly\alpha}}\right)^2.
\end{equation}

%-----------------------------------------------------------------------------
\subsection{Gaussian Detuning Scaling}
%-----------------------------------------------------------------------------

For a symmetric line profile with thermal width $\sigma$ and small detuning 
$\delta \ll \sigma$, a Taylor expansion of the Gaussian gives:
\begin{equation}
A = \frac{\Delta I}{I} \propto \left(\frac{\delta}{\sigma}\right)^2
\end{equation}
to leading order (the linear term vanishes by symmetry). This scaling follows 
from the sensitivity of resonant absorption/scattering to wavelength mismatch.

%-----------------------------------------------------------------------------
\subsection{The Double-Transit Mechanism}
%-----------------------------------------------------------------------------

\paragraph{Physical picture.}
Chromospheric Ly-$\alpha$ photons are resonantly scattered by coronal hydrogen 
atoms before reaching the observer. In DFD, this involves two passages through 
the refractive corona:
\begin{enumerate}[nosep]
\item \textbf{Incoming leg:} Chromosphere $\to$ scattering site in corona
\item \textbf{Outgoing leg:} Scattering site $\to$ observer
\end{enumerate}

Locally-produced O~VI emission involves only one passage:
\begin{enumerate}[nosep]
\item \textbf{Outgoing leg:} Emission site in corona $\to$ observer
\end{enumerate}

\paragraph{Detuning accumulation.}
Let $\delta_{\rm in}$ be the detuning accumulated on the incoming leg and 
$\delta_{\rm out}$ be the detuning on the outgoing leg. The double-transit 
hypothesis asserts:
\begin{align}
\delta_{\rm Ly\alpha} &= \delta_{\rm in} + \delta_{\rm out} \approx 2\delta_0, \\
\delta_{\rm OVI} &= \delta_{\rm out} \approx \delta_0,
\end{align}
where $\delta_0$ is a characteristic detuning per leg.

\paragraph{Resulting enhancement.}
With $A \propto (\delta/\sigma)^2$:
\begin{equation}
\Gamma = \frac{(\delta_{\rm Ly\alpha})^2}{(\delta_{\rm OVI})^2} = \frac{(2\delta_0)^2}{(\delta_0)^2} = 4.
\end{equation}

%-----------------------------------------------------------------------------
\subsection{The Conservative-Field Consistency Check}
%-----------------------------------------------------------------------------

A careful reader may object: if the DFD shift is governed by a conservative 
scalar field $\psi$, then accumulated phase/wavelength changes depend only on 
\emph{endpoints}:
\begin{equation}
\int_{\rm path} \nabla\psi \cdot d\ell = \psi(\text{end}) - \psi(\text{start}),
\end{equation}
independent of the geometric path length. In that case, ``two passes through 
the same region doubles the shift'' is not automatic.

\paragraph{Resolution.}
The double-transit effect does not require path-length dependence of $\psi$. 
Rather, it arises from the \emph{measurement geometry}: the UVCS asymmetry 
statistic compares different sightlines (east vs.\ west limb), and the relevant 
quantity is the \emph{differential} detuning between directions.

For scattered Ly-$\alpha$:
\begin{itemize}[nosep]
\item The incoming photon samples the $\psi$ gradient from chromosphere to 
scattering site
\item The outgoing photon samples the $\psi$ gradient from scattering site to 
observer
\item Both gradients contribute to the E-W asymmetry
\end{itemize}

For locally-emitted O~VI:
\begin{itemize}[nosep]
\item Only the outgoing leg contributes
\end{itemize}

The key assumption is: \emph{the detuning relevant for the asymmetry $A$ receives 
additive contributions from both legs for resonantly scattered Ly-$\alpha$, while 
the O~VI statistic samples only one leg.}

This assumption should be verified against the explicit UVCS measurement 
definition, which is why we present $\Gamma$ as a measured quantity rather 
than an assertion.

%-----------------------------------------------------------------------------
\subsection{Observational Constraint on $\Gamma$}
%-----------------------------------------------------------------------------

From the UVCS data:
\begin{align}
R_{\rm obs} &= 39.2 \pm 8.2, \\
\left(\frac{\sigma_{\rm OVI}}{\sigma_{\rm Ly\alpha}}\right)^2 &= 9.0.
\end{align}

Direct inversion gives:
\begin{equation}
\boxed{\Gamma_{\rm obs} = \frac{R_{\rm obs}}{9} = 4.4 \pm 0.9}
\end{equation}

This is consistent with the double-transit prediction $\Gamma = 4$ at $0.4\sigma$, 
and inconsistent with the standard physics prediction $\Gamma = 1$ at $3.7\sigma$.

%-----------------------------------------------------------------------------
\subsection{Falsifiable Predictions}
%-----------------------------------------------------------------------------

The $\Gamma = 4$ hypothesis makes crisp empirical predictions that can be tested 
with existing or future data:

\paragraph{1. Scattered vs.\ local lines.}
Other lines dominated by resonant scattering should share $\Gamma \approx 4$:
\begin{itemize}[nosep]
\item H-$\alpha$ (if observable in scattered component)
\item He~II 304~\AA\ (scattered transition-region emission)
\end{itemize}
Purely collisional coronal lines should show $\Gamma \approx 1$:
\begin{itemize}[nosep]
\item Fe~XII 195~\AA
\item Fe~XIV 211~\AA
\item Mg~X 625~\AA
\end{itemize}

\paragraph{2. Geometry dependence.}
If $\Gamma$ arises from two-leg sampling, it should vary with viewing geometry:
\begin{itemize}[nosep]
\item Limb observations: Maximum scattering geometry, largest $\Gamma$
\item Disk center: Minimal scattering toward observer, reduced $\Gamma$
\end{itemize}
The predicted variation can be calculated from the scattering phase function.

\paragraph{3. Hybrid lines.}
Lines with mixed collisional + scattered contributions should show intermediate 
$\Gamma$ values, weighted by the fractional contributions.

\paragraph{4. Solar cycle variation.}
If coronal conditions affect the relative contributions of scattered vs.\ local 
emission, $\Gamma$ may vary with solar activity level.

%-----------------------------------------------------------------------------
\subsection{Summary}
%-----------------------------------------------------------------------------

The UVCS asymmetry ratio provides a clean test of DFD's refractive mechanism:

\begin{center}
\begin{tabular}{lcc}
\toprule
\textbf{Model} & \textbf{Predicted $\Gamma$} & \textbf{Status} \\
\midrule
Standard physics & 1 & Excluded at $3.7\sigma$ \\
DFD (double-transit) & 4 & Consistent at $0.4\sigma$ \\
\midrule
\textbf{Observed} & $4.4 \pm 0.9$ & --- \\
\bottomrule
\end{tabular}
\end{center}

The double-transit derivation converts the enhancement factor from an assertion 
into a \emph{measurable prediction} with explicit falsifiers. Future observations 
of additional line species and geometries can definitively confirm or refute 
$\Gamma = 4$.
